\problem Два графа \textit{изоморфны}, если вершины одного можно занумеровать вершинами другого так, что рёбра перейдут в рёбра. Сколько существует попарно неизоморфных деревьев на шести вершинах? Нарисуйте их.

\problem Вершину степени 1 называют \textit{висячей}, или \textit{листом}. Докажите, что в дереве, у которого больше одной вершины, есть висячая вершина, и что таких вершин хотя бы две.

\problem Докажите, что в дереве с $n$ вершинами ровно $n-1$ ребро.

\problem Сколько рёбер в \textit{лесу} — графе без циклов — с $n$ вершинами и $k$ компонентами связности?

\problem Докажите, что из любого связного графа можно выкинуть часть рёбер так, чтобы получилось дерево с теми же вершинами. Такое дерево называют \textit{остовом}, или \textit{остовным деревом}, графа.

\problem В графе $n$ вершин и $m$ рёбер. Докажите, что компонент связности в нём не меньше $n-m$.

\problem Пусть граф связен и имеет $n$ вершин. Докажите, что следующие свойства равносильны: $\bullet$ рёбер ровно $n-1$; $\bullet$ при удалении любого ребра граф перестаёт быть связным; $\bullet$ в графе нет циклов; $\bullet$ любые две вершины соединены ровно одним путём.\\
Таким образом, любое из этих свойств можно взять за определение дерева.

\problem Даны $n$ различных по весу камней и весы, позволяющие сравнить любые два из них. Какое наименьшее число взвешиваний нужно, чтобы наверняка найти самый тяжёлый камень?

\problem Сколько рёбер может быть в несвязном графе с $n$ вершинами?

\problem Докажите, что из любого связного графа, в котором больше одной вершины, можно удалить вершину вместе со всеми выходящими из неё рёбрами так, что граф останется связным.

\problem В связном графе 10 вершин и 19 рёбер. Докажите, что в нём найдётся цикл, после уничтожения всех рёбер которого граф не потеряет связность.

\problem В стране 100 городов, некоторые из которых соединены авиалиниями. Известно, что из любого города можно долететь до любого другого, возможно с пересадками. Докажите, что можно побывать в каждом городе, совершив не более 198 перелётов.

\problem Дано дерево с $n$ вершинами. Сколькими способами можно раскрасить его вершины в $k$ цветов так, чтобы концы любого ребра были разного цвета?

\problem Выпуклый $n$-угольник разрезали на треугольники непересекающимися диагоналями.\\
а) Докажите, что треугольников получилось на один больше, чем проведённых диагоналей.\\
б) Докажите, что хотя бы у двух треугольников по две стороны лежат на границе многоугольника.\\
в) Сколькими способами можно так разрезать многоугольник? Вершины занумерованы: разрезания, отличающиеся поворотом, считаются различными.

\problem Хозяйка испекла пирог и знает, что придут либо ровно $p$, либо ровно $q$ гостей, причём $p$ и $q$ взаимно просты. На какое наименьшее число не обязательно равных кусков можно разрезать пирог, чтобы в любом случае разделить его между гостями поровну?
